\sscp{*kusor `cuscus'}{14-03-02}
The Papuan \tsc{North Halmahera} languages have forms for ‘cuscus’
which are reflexes of *kusor.
Known reflexes of these forms are given in \trf{14-02}.
Stress in Sahu is indicated by an acute accent.

\begin{table}[h]
	\centering\tcp{Reflexes of *kusor cuscus in \tsc{North Halmahera} languages}{14-02}
	\st{6pt}\begin{tabular}{l|lll}\lsptoprule
Proto-NH	&	*kusor	&	Gloss	&	Source \\ \midrule
Sahu	&	ʔúsoro	&	cuscus	&	\citet[131]{VisserVoorhoeve1987} \\
Ternate	&	kuso	&	cuscus	&	\citet[269]{Schapper2011} \\
Tidore	&	kuso	&	cuscus	&	\citet[269]{Schapper2011} \\
Galela	&	kuso	&	cuscus	&	\citet[223]{vanBaarda1895} \\
Tobelo	&	kuho	&	cuscus	&	\citet[212]{Hueting1908}\\
Pagu	&	kuso	&	cuscus	&	Dalan Mehuli Perangin Angin \\
			&						&					&	(pers. comm. May 2023)  \\
West Makian	&	kusu	&	cuscus	&	\citet[81]{Collins1982b} \\
		\lspbottomrule
	\end{tabular}
\end{table}

On the basis of these forms,
\citet[269]{Schapper2011} reconstructs Proto-North Halmahera *kusoro.
However, more detailed investigation shows *kusor is a more appropriate reconstruction.
The Sahu form with three syllables adds a copy vowel after a
historical consonant-final form, while the other cognates have lost the historical word-final consonant.
Extrametrical copy vowels (also called echo vowels in some literature
of \tsc{North Halmahera} languages) are quite common in these languages.

Examination of a wider range of data shows that most Sahu trisyllables
(including \it{ʔúsoro}, ‘cuscus’) are of the shape CV́\sub{1}CV\sub{2}CV\sub{2}; that
is the penultimate and final vowels are identical,
and stress falls on the antepenultimate syllable.
Comparison with other \tsc{North Halmahera} languages reveals
that such trisyllables frequently correspond to a vowel-final
disyllable in Galela, Ternate, Tidore, and West Makian, but a consonant-final disyllable in Pagu.
The correspondences of Sahu CV́\sub{1}CV\sub{2}CV\sub{2} trisyllables
with cognate forms in Tobelo are either disyllables, or CV́\sub{1}CV\sub{2}CV\sub{2} trisyllables.
Examples are given in \trf{14-03},
along with tentative Proto-North Halmahera reconstructions.
Most words selected for \trf{14-03} have a final liquid /r/ or
/l/ for maximum comparability with *kusor ‘cuscus’.

\begin{table}[h]\centering\tcp{Word-final consonants in \tsc{North Halmahera}}{14-03}
\begin{adjustbox}{width=\textwidth}
	\begin{threeparttable}\st{2pt}
	\begin{tabular}{l|llllllll}\lsptoprule
*gloss	&	tears (n.)	&	liver	&	quiet	&	neck	&	belly	&	palm	&	hand	&	fish\\
Proto-NH	&	*koŋor	&	*gater	&	*ogor\su{‡}	&	*ʧomal	&	*pokol	&	*bobol\su{Δ}	&	*giam	&	*ɲawok\\	\midrule
Pagu	&	koŋol	&	gatel	&	ogol	&	tomal	&	pokol	&	bobol	&	giam	&	naok\\
Sahu	&	ʔóŋoro	&	kátere	&	ógoro	&	ʧámala	&	póʔolo	&	bóbolo	&	gíama	&	ɲáoʔo\\
Galela	&	koŋo	&	gate\su{†}	&	ogo	&	tolo	&	poko	&	boboro	&	gia	&	nawo\\
Tobelo	&	koŋo	&	gate	&	ogoro	&	tomara	&	pokoro	&	boboro	&	giama	&	nawoko\\
Ternate	&	oŋo	&	gate	&	ogo	&	ʧama	&	\cg{(oru)}	&	bobo	&	gia	&	ɲau\\
Tidore	&	oŋo	&	gate	&	googo	&	\cg{(saka)}	&	\cg{(oru)}	&		&	gia	&	ɲao\\
W. Makian	&		&	\cg{(amo)}	&		&	\cg{(ru)}	&		&		&	ia	&	yao\\
		\lspbottomrule
	\end{tabular}
		\begin{tablenotes}\tnsize
			\item[†]	{Galela \it{gate} liver is from \citet{Wada1980}.
								We were unable to locate it in \citet{vanBaarda1895}.}
			\item[‡]	{The Pagu, Galela and Tobelo reflexes of *ogor are given as `cease of sth.
								like wind, rain, waves'.
								The Sahu, Ternate and Tidore reflexes are `quiet, still'.}
			\item[Δ]	{Pagu and Ternate reflexes of *bobol are `lontar palm'.
								Other forms are `nipa palm'.}
		\end{tablenotes}
	\end{threeparttable}
\end{adjustbox}
\end{table}

Such words with copy vowels contrast with Sahu trisyllables which have penultimate stress.
Such words can have different penultimate
and final vowels (CV\sub{1}CV́\sub{2}CV\sub{3}).
These trisyllables correspond to trisyllables in other North Halmahera languages.
Examples are given in \trf{14-04}.

\begin{table}[h]\centering\tcp{Proto-North Halmahera trisyllables}{14-04}
\begin{adjustbox}{width=\textwidth}
	\begin{threeparttable}\st{2pt}
	\begin{tabular}{l|lllllll}\lsptoprule
gloss	&	yellow, gold	&	weaver ant	&	round	&	middle	&	bind	&	ginger	&	eyebrow\\
PNH	&	*[k/g]urati	&	*gunaŋe\su{†}	&	*[b/p]olulu	&	*korona	&	*podiku	&		&	*goroputu\\	\midrule
Pagu	&	kulati	&	gunane	&	palulun	&	kolona	&	piliku	&		&	goloutu\\
Sahu	&	guráʧi	&	gunáne	&	ɓulúlu	&	ʔolóna	&	piríʔu	&	galáʔa	&	gorowútu\\
Galela	&	kurati	&	gunene	&	pululu, bululu	&	konora\su{‡}	&	piliku	&		&	\\
Tobelo	&	kurati	&		&	bululu	&	korona	&	\cg{(ʎikutu)}	&		&	gorohutu\\
Ternate	&	kuraʧi	&	gunaŋe	&	polulu	&	konora	&	podiku	&	goraka	&	\\
Tidore	&	kuraʧi, guraʧi	&		&	folulu	&	konora	&		&	goraka	&	\\
W. Makian	&	burei	&		&		&		&		&	buree	&	\\
		\lspbottomrule
	\end{tabular}
		\begin{tablenotes}\tnsize
			\item[†]	{Pagu \it{gunane} is given as `red forest ant',
								Sahu \it{gunáne} as `k.o.
								ant living on trees
								and making nests of folded leaves',
								Galela \it{gunene} as `chicken louse,
								scabies mite; and other invisible tiny itchy insects',
								Ternate \it{gunaŋe} as `weaver ant', \it{Oecophylla smaragdina}.}
			\item[‡]	{Galela \it{konora} means `between'.
								It is identified by \citet[215]{vanBaarda1895} as a loan from Ternate.}
		\end{tablenotes}
	\end{threeparttable}
\end{adjustbox}
\end{table}

The best analysis of the data 
in \trf{14-03} is probably that they are
reflexes of earlier consonant final forms.
Sahu has undergone epenthesis of a final copy vowel to create
a vowel-final form, while most other languages have lost final consonants; *C > {\0} /{\gap}{\#}.
Tobelo has undergone both processes with any conditioning environment currently unclear.\footnote{
		A process of copy vowel epenthesis has also affected certain consonant
		final loans in some \tsc{North Halmahera} languages.
		Two examples from Sahu are Malay \it{campur} [ʧampur] > Sahu \it{ʧápuru} ‘mix’,
		and Malay \it{durian} > Sahu \it{duríana} ‘durian’.}
Pagu is conservative in retaining word-final consonants unchanged.
While Pagu
retains word-final consonants, \citet[78]{Wimbish1992} describes a process of vowel epenthesis:
“There is a tendency in some dialects [of Pagu] to add a copy
vowel [{\ldots}] to many consonant-final roots.
This may be resultant from influence of related languages which
do not allow consonant-final roots [{\ldots}] Although speakers
frequently drop the added vowel when asked for the root word,
in natural speech it is usually present.
The addition of this vowel results in examples with stress which
is not penultimate.” 
Given all this data,
the forms meaning ‘cuscus’ in the \tsc{North Halmahera} languages
presented in \trf{14-02} are best accounted for by Proto-North Halmahera *kusor ‘cuscus’.
While Pagu has contrastive liquids /l/
and /r/, of these only /l/ is allowed word finally \citep[78]{Wimbish1992}.
The Pagu form \it{kuso} without a final consonant is thus unexpected,
as the regular form based on other correspondences would be unattested **kusol.
This might point to distribution through contact.